\section{Ejecución del bloque II: P05--P08}
\label{sec:bloque2}
\textbf{Fecha: 18 de octubre de 2026.} Se implementaron propiedad del estado,
rutas por modalidad, procedencia de restricciones y un fallback conservador.
Se mantienen pendientes las mediciones con modelos reales y anotaciones humanas.
Este bloque cambia algunos estados de la API y reduce deliberadamente la
cobertura automática de interpretaciones inciertas; no es una mejora universal.

\subsection{P05: estado serializable y propietario único de exclusiones}
\textbf{Antes.} Juez y grafo podían escribir el pool. La idempotencia evitaba
duplicados, pero el estado contenía un objeto de pool y una caché, inadecuados
como contrato de persistencia.

\textbf{Ahora.} El juez devuelve decisiones sin mutar exclusiones. El único
punto de escritura del flujo es el nodo de evaluación del orquestador.
El estado contiene listas y diccionarios serializables: decisiones, exclusiones,
versión de plan, candidatos y evidencia. La caché de recuperación queda fuera
del estado, ligada a la instancia de grafo creada para una solicitud.
El antiguo argumento de pool se admite por compatibilidad, pero se ignora;
el tipo auxiliar no participa en el flujo actual.

Cada exclusión registra candidato, motivo, iteración, alcance y versión del
plan. Los datos inexistentes se excluyen para toda la solicitud; las decisiones
dependientes de las consultas se vinculan al plan evaluado. Un rechazo histórico
no reaparece como alternativa sin evaluar: solo el candidato seleccionado puede
volver después de reevaluarse con un plan nuevo. Se conserva la decisión previa
en el historial.

\textbf{Evidencia.} Pruebas de estado convertible a JSON sin objetos opacos,
decisión pura del juez, escritura única por rechazo y solicitudes concurrentes
con estados separados. \textbf{Ventaja.} Se puede auditar quién decidió qué y
con qué plan; desaparece la doble responsabilidad de mutación.
\textbf{Límite.} No se añadió almacenamiento durable ni reanudación. La caché
fuera del estado no se guarda al serializar. Reutilizar el grafo o añadir
checkpoints necesita el trabajo de P16/P18.

\subsection{P06: modalidades e instrumentales sin aprobación acústica ficticia}
\textbf{Antes.} El flujo exigía letras a cualquier candidato e invocaba al juez
incluso para consultas acústicas. Una aprobación lírica podía confundirse con
una aprobación multimodal.

\textbf{Ahora.} Una consulta solo acústica omite el juez lírico, admite candidatos
sin letra y termina como \code{retrieved}: recuperado, con audio
\code{not_evaluated}. Una consulta solo lírica puede terminar como
\code{approved}. Si la letra se aprueba y también se pidió audio, la terminación
es \code{partially_validated}; el sonido sigue sin evaluación independiente.

Cada resultado informa \code{validation.lyrics} y \code{validation.audio}, con
estados de aprobado, rechazado, error, no solicitado o no evaluado. La interfaz
muestra esos estados. Una ausencia de letra en una solicitud que sí la requiere
se registra como falta de evidencia, no como rechazo de un LLM que no fue llamado.
\code{judge_iterations} cuenta invocaciones reales al juez; el contador de
selecciones queda en \code{candidate_iterations}.

\textbf{Evidencia.} Instrumentales sin letra, omisión de llamadas al juez,
validación parcial en consultas mixtas y respuestas de Flask real.
\textbf{Ventaja.} Evita descartar instrumentales por una condición ajena a su
modalidad y declara el alcance de la evidencia.
\textbf{Límite.} CLAP sigue proporcionando recuperación, no certificación de
instrumentos, voces o energía. La calidad de la ruta acústica no se ha medido.
Las restricciones líricas negativas también requieren evaluación aunque no haya un tema positivo lírico. El presupuesto todavía cuenta selecciones; su separación completa sigue en P11.

\subsection{P07: restricciones con origen y reparación del plan}
\textbf{Antes.} Todos los negativos quedaban inmóviles, aunque fueran una
inferencia equivocada. Tampoco podía activarse una modalidad omitida por el router.

\textbf{Ahora.} \code{intent.py} prepara restricciones con ID, modalidad,
polaridad, texto, origen y cita. \code{user_explicit} exige una cita literal del
usuario; una negativa debe incluir una negación reconocida. Lo no fundamentado
se registra como \code{model_hypothesis}. Las restricciones explícitas capturadas
se conservan; las hipótesis pueden corregirse o retirarse.

Una modalidad puede activarse o desactivarse para corregir la intención si se
aporta una cita y diagnóstico de esa corrección. Cambiar audio siempre requiere
evidencia textual con una señal de modalidad soportada: un rechazo lírico por
sí solo no justifica describir como incorrecto el sonido de una canción no oída.
Un cambio de plan se versiona y sus decisiones se reevaluarán bajo esa versión.
Las modalidades no seleccionadas permanecen intactas.

Se rechaza un plan que omita anclar una negación visible del usuario. Este
control usa señales léxicas, no una comprensión completa del lenguaje. El prompt
del router ahora pide procedencia y listas vacías para modalidades no solicitadas.
Se probó el formateo con la plantilla LangChain real para que los ejemplos JSON
no introduzcan variables de plantilla inexistentes.

\textbf{Evidencia.} Pruebas de eliminación de hipótesis, conservación de
restricciones explícitas, citas inventadas o sin negación, corrección de
modalidades y reevaluación después de cambiar de plan.
\textbf{Ventaja.} Separa decisiones revisables de restricciones que el usuario
declaró y evita perderlas silenciosamente.
\textbf{Límite importante.} Verificar una cita no demuestra que su traducción
inglesa sea semánticamente fiel. Un modelo todavía puede inferir demasiado al
etiquetar una restricción como explícita. Se registra E13 para calibrarlo.

\subsection{P08: fallback limitado, trazable y con abstención}
\textbf{Antes.} El respaldo copiaba la consulta cruda en ambas modalidades,
con negativos vacíos, sin garantía de idioma o intención.

\textbf{Ahora.} Puede reutilizarse un plan válido únicamente si está asociado
a la misma consulta exacta. La función permite esa recuperación, pero no guarda
un historial global de consultas ni se conecta todavía a una persistencia.
Cuando no hay plan previo, una gramática local acotada interpreta frases de
sonido soportadas en español e inglés, conserva las negaciones reconocidas y
registra sus citas. Por ejemplo, piano lento sin voces produce consultas
acústicas positivas y una negativa de voces, con la rama lírica vacía.

Si quedan palabras o alcances de negación sin interpretar, se devuelve
\code{clarification_needed}, listas de búsqueda vacías y una advertencia de
reformulación. No se llama a workers ni al juez. No se promete traducción
general, interpretación lírica offline ni extracción arbitraria de negativos.
La respuesta HTTP es 200 con ese estado: el JSON es válido, pero no se ejecutó
una búsqueda de significado incierto.

\textbf{Evidencia.} Frases bilingües soportadas, restricciones negativas,
negaciones ambiguas, planes previos de otra consulta y fallo del router.
\textbf{Ventaja.} La degradación no presenta como fiel una consulta que perdió
restricciones. \textbf{Coste.} Hay consultas para las que el respaldo devuelve
menos resultados o ninguno y exige reformular; la cobertura y satisfacción
necesitan medición antes de ampliar reglas.

\subsection{Resultados de control y regresiones observadas}
La suite pasó 72 pruebas en el entorno del bloque I, con Flask, FAISS y
LangGraph reales y proveedores/modelos simulados. Se actualizaron expectativas
anteriores para los estados nuevos y el juez sin efectos laterales; no se
eliminaron los casos de regresión originales.

Se ejecutaron 36 corridas simuladas sin refinamiento y 36 con refinamiento.
En cada serie se observaron 12 aprobaciones líricas, 6 validaciones parciales,
12 recuperaciones acústicas sin validación independiente y 6 solicitudes de
reformulación. Estas últimas corresponden a dos consultas repetidas tres veces:
\code{negative_romance} y \code{instrumental_study}. Sus planes sintéticos
anteriores no incluyen la procedencia requerida para las negaciones y la
gramática limitada no cubre toda la petición. El fixture original se mantuvo
para que ese cambio de cobertura sea visible.

No hubo errores del flujo. Se registran por separado reconsideraciones históricas
y reapariciones de excluidos sin reevaluación; las segundas deben permanecer en
cero. Jaccard solo se calcula cuando hay resultados; una respuesta sin ranking
no recibe un 1 ficticio. Las métricas musicales continúan nulas. La ejecución
real se mantuvo bloqueada sin llamar al proveedor.

Esto demuestra controles adicionales, no mejor relevancia musical. La reducción
de cobertura automática es una contrapartida explícita; E02/P12 deben evaluar
si conviene ampliar el respaldo o cambiar la estrategia con evidencia real.
Los informes del bloque I se conservaron; los actuales están en
\code{validation/block2/} y llevan el identificador \code{block2-pre-P09}.

\subsection{Compatibilidad y trabajo que no se cierra}
La API añade \code{retrieved}, \code{partially_validated} y
\code{clarification_needed}; clientes que esperaban solo aprobación/error/límite
deben interpretar estos estados. \code{validation_scope} puede ser
\code{none}; \code{constraints}, \code{plan_version} y validación por resultado
permiten inspeccionar la evidencia. El argumento legado de pool ya no aplica
exclusiones por sí mismo. Los locks no cambiaron y no se instalaron dependencias.

P05--P08 tienen implementación y verificación de control. Se cubren la propiedad
única de E08, las rutas de E01 y los controles de E02/E03. Persistencia, evidencia
acústica independiente, fidelidad semántica y evaluación real siguen abiertas.
Los siguientes puntos de implementación son P09--P12, no la ampliación de la UI
o ingesta. La recuperación parcial ante fallos del refinador sigue pendiente.
